\documentclass[10pt,a4paper]{amsart}
\usepackage{amsmath,amssymb,amsthm,graphicx}
\usepackage{verbatim}
\newenvironment{reference}{\comment}{\endcomment}
%\newenvironment{reference}{}{}
\newenvironment{slogan}{\comment}{\endcomment}
\newenvironment{history}{\comment}{\endcomment}

% For commutative diagrams we use Xy-pic
\usepackage[all]{xy}

% We use 2cell for 2-commutative diagrams.
\xyoption{2cell}
% Theorem environments.
%
\theoremstyle{plain}
\newtheorem{theorem}[subsection]{Theorem}
\newtheorem{proposition}[subsection]{Proposition}
\newtheorem{lemma}[subsection]{Lemma}

\theoremstyle{definition}
\newtheorem{definition}[subsection]{Definition}
\newtheorem{example}[subsection]{Example}
\newtheorem{exercise}[subsection]{Exercise}
\newtheorem{situation}[subsection]{Situation}

\theoremstyle{remark}
\newtheorem{remark}[subsection]{Remark}
\newtheorem{remarks}[subsection]{Remarks}

\numberwithin{equation}{subsection}

% Macros
%
\def\lim{\mathop{\mathrm{lim}}\nolimits}
\def\colim{\mathop{\mathrm{colim}}\nolimits}
\def\Spec{\mathop{\mathrm{Spec}}}
\def\Hom{\mathop{\mathrm{Hom}}\nolimits}
\def\Ext{\mathop{\mathrm{Ext}}\nolimits}
\def\SheafHom{\mathop{\mathcal{H}\!\mathit{om}}\nolimits}
\def\SheafExt{\mathop{\mathcal{E}\!\mathit{xt}}\nolimits}
\def\Sch{\mathit{Sch}}
\def\Mor{\mathop{\mathrm{Mor}}\nolimits}
\def\Ob{\mathop{\mathrm{Ob}}\nolimits}
\def\Sh{\mathop{\mathit{Sh}}\nolimits}
\def\NL{\mathop{N\!L}\nolimits}
\def\CH{\mathop{\mathrm{CH}}\nolimits}
\def\proetale{{pro\text{-}\acute{e}tale}}
\def\etale{{\acute{e}tale}}
\def\QCoh{\mathit{QCoh}}
\def\Ker{\mathop{\mathrm{Ker}}}
\def\Im{\mathop{\mathrm{Im}}}
\def\Coker{\mathop{\mathrm{Coker}}}
\def\Coim{\mathop{\mathrm{Coim}}}

% Boxtimes
%
\DeclareMathSymbol{\boxtimes}{\mathbin}{AMSa}{"02}

%
% Macros for moduli stacks/spaces
%
\def\QCohstack{\mathcal{QC}\!\mathit{oh}}
\def\Cohstack{\mathcal{C}\!\mathit{oh}}
\def\Spacesstack{\mathcal{S}\!\mathit{paces}}
\def\Quotfunctor{\mathrm{Quot}}
\def\Hilbfunctor{\mathrm{Hilb}}
\def\Curvesstack{\mathcal{C}\!\mathit{urves}}
\def\Polarizedstack{\mathcal{P}\!\mathit{olarized}}
\def\Complexesstack{\mathcal{C}\!\mathit{omplexes}}
% \Pic is the operator that assigns to X its picard group, usage \Pic(X)
% \Picardstack_{X/B} denotes the Picard stack of X over B
% \Picardfunctor_{X/B} denotes the Picard functor of X over B
\def\Pic{\mathop{\mathrm{Pic}}\nolimits}
\def\Picardstack{\mathcal{P}\!\mathit{ic}}
\def\Picardfunctor{\mathrm{Pic}}
\def\Deformationcategory{\mathcal{D}\!\mathit{ef}}
\usepackage{cleveref}
\crefname{originalproposition}{original proposition}{original propositions}
\Crefname{originalproposition}{Original proposition}{Original propositions}
\makeatletter
\newlabel{algebra-proposition-fp-tensor}{{059K}{source}{}{sourcepointer-1}{}}
\newlabel{algebra-proposition-fp-tensor@cref}{{[originalproposition][0][]059K}{[1][0][]source}{}{}{}}
\makeatother
\crefname{originalreference}{original reference}{original references}
\Crefname{originalreference}{Original reference}{Original references}
\makeatletter
\newlabel{lemma-countable-fg-tensor}{{0AL9}{source}{}{sourcepointer-2}{}}
\newlabel{lemma-countable-fg-tensor@cref}{{[originalreference][0][]0AL9}{[1][0][]source}{}{}{}}
\makeatother
\title{For a countable module over a countable ring, finite presentation is equivalent to the canonical map M tensor R\textasciicircum{}N -> M\textasciicircum{}N being an isomorphism}
\author{The Stacks Project Authors}
\date{stacks-project:examples.tex:local-snapshot-20261012; GNU FDL 1.2 or later}
\begin{document}
\maketitle
\noindent\textit{Editorial note.} All delivered mathematical text is the author's own, copied from the pinned Stacks Project chapter file examples.tex; no statement or proof step was rewritten, reordered or removed.\par\medskip
\noindent\textit{Editorial note.} The counted proof is the author's complete argument: The kernel-countability step is proved by 'arguing as in the proof of Algebra, Proposition algebra-proposition-fp-tensor' - the same enumeration argument as the finite-generation case applied to K; recorded as mechanical analogy, content verbatim.\par\medskip
\noindent\textit{Editorial note.} The label lemma-countable-fp-tensor names the counted statement, and the proof environment immediately following it in the source is delivered verbatim as its proof.\par\medskip
\noindent\textit{Editorial note.} The counted statement is carried by the author's own lemma environment in the chapter source.\par\medskip
\noindent\textit{Allowed standard prerequisites (editorial).}\begin{itemize}
\item The canonical map M \textbackslash{}otimes\_R R\textasciicircum{}N -> M\textasciicircum{}N and its finitely generated counterpart (Algebra, proposition-fp-tensor); countable modules and countable rings
\item Same-chapter lemma-countable-fg-tensor (L48287 cluster): for countable M, finite generation iff the map is surjective
\end{itemize}\medskip
\ifdefined\hypertarget\hypertarget{sourcepointer-1}{}\fi
\noindent\textit{Source reference.} Original proposition 059K: sources/stacks-chow/tags.txt, source lines 6818--6818. Prerequisite pointer invoked at 988; it names a standard fact of a sibling chapter that the counted argument uses, and every step of the counted proof built on it is retained verbatim in the delivered body.\par\medskip
\ifdefined\hypertarget\hypertarget{sourcepointer-2}{}\fi
\noindent\textit{Source reference.} Original reference 0AL9: sources/stacks-examples/pinned-original-source.tex, source lines 963--963. Prerequisite pointer invoked at 989; it names a standard fact of the same chapter that the counted argument uses, and every step of the counted proof built on it is retained verbatim in the delivered body.\par\medskip
\begin{lemma}
\label{lemma-countable-fp-tensor}
Let $R$ be a countable ring. Let $M$ be a countable $R$-module. Then $M$
is finitely presented if and only if the canonical map
$M \otimes_R R^\mathbf{N} \to M^\mathbf{N}$ is an isomorphism.
\end{lemma}

\begin{proof}
If $M$ is a finitely presented module, then the map is an isomorphism
by Algebra, Proposition \ref{algebra-proposition-fp-tensor}. Conversely,
assume the map is an isomorphism. By Lemma \ref{lemma-countable-fg-tensor}
the module $M$ is finite. Choose a surjection $R^{\oplus m} \to M$ with
kernel $K$. Then $K$ is countable as a submodule of $R^{\oplus m}$.
Arguing as in the proof of Algebra, Proposition
\ref{algebra-proposition-fp-tensor} we see that
$K \otimes_R R^\mathbf{N} \to K^\mathbf{N}$ is surjective.
Hence we conclude that $K$ is a finite $R$-module by
Lemma \ref{lemma-countable-fg-tensor}.
Thus $M$ is finitely presented.
\end{proof}

\end{document}
